%% PLANTILLA PARA ARTICULOS
\documentclass[10pt,twoside,spanish,a4paper]{article}
%%%%%%%%%%%%%%%%%%%%%%%%%%%%% no tocar
% Alfredo Cuesta. 01/01/2013. acuestai@ucm.es

%si se quiere usar ARIAL hay que poner estas 2 lineas:
%\usepackage[scaled]{uarial} %DESCOMENTAR
\renewcommand*\familydefault{\sfdefault}
%

\usepackage[compact]{titlesec}
 \titlespacing\section{0pt}{12pt plus 0pt minus 0pt}{0pt plus 0pt minus 0pt}
 \titlespacing\subsection{0pt}{5pt plus 0pt minus 0pt}{0pt plus 0pt minus 0pt}
 \titlespacing\subsubsection{0pt}{0pt plus 0pt minus 0pt}{0pt plus 0pt minus 0pt}

\usepackage[utf8]{inputenc}
\usepackage{
amssymb,
graphicx,
color,
url
%wrapfig %Inclusión de gr·ficos al lado de texto
}
\definecolor{olivegreen}{rgb}{0.0,0.5,0.0}

% \usepackage{listings}
% \lstset{language=Matlab,%
%     basicstyle=\footnotesize,
%     keywordstyle=\color{blue},%
%     commentstyle=\color{olivegreen}\textit,%
%     emph=[1]{ans},emphstyle=[1]\color{black}, %some words to emphasise
%     emph=[2]{resample,timeseries,syms},emphstyle=[2]\color{blue}, %some words to emphasise
% }

\linespread{1.3}%separacion de lineas
\parindent=0.0in%indentancion de parrafo
\parskip=6pt%separacion entre el parrafo siguiente
\setlength{\topmargin}{-0.5in}
\setlength{\oddsidemargin}{0in}
\setlength{\evensidemargin}{0in}
\setlength{\textwidth}{16cm}%{4.72in}
\setlength{\textheight}{22cm}%{7.08in}

\usepackage{babel}
\addto\shorthandsspanish{\spanishdeactivate{~<>}}
% correct bad hyphenation here
\hyphenation{nu-es-tro co-no-ci-mien-to par-ti-cu-lar ca-rac-te-ris-ti-cas ni-vel se-gui-mien-to que-re-mos ca-ta-lo-go}
\date{}


% --- customization ---

\renewcommand\labelitemi{$\bullet$}
\let\olditemize\itemize
\renewcommand{\itemize}{
  \olditemize
  \vspace{-0pt}
  \setlength{\itemsep}{6pt}
  \setlength{\parskip}{0pt}
  \setlength{\parsep}{0pt}
  %\setlength{\topsep}{-15pt}
}

\let\oldenumerate\enumerate
\renewcommand{\enumerate}{
  \oldenumerate
  \vspace{-2pt}
  \setlength{\itemsep}{1pt}
  \setlength{\parskip}{0pt}
  \setlength{\parsep}{0pt}
  %\setlength{\topsep}{-15pt}
}

\usepackage{enumitem}
\usepackage{fancyhdr}
\usepackage{multirow}
\usepackage{colortbl}
\usepackage{amsmath}
% \usepackage[framed]{mcode}
\usepackage{matlab-prettifier}
\lstset{literate={ñ}{{\~n}}1 {á}{{\'a}}1 {é}{{\'e}}1 {í}{{\'i}}1 {ó}{{\'o}}1 {ú}{{\'u}}1}

\begin{document}

\renewcommand{\tablename}{Tabla}
\renewcommand{\figurename}{Figura}
\title{
  \vspace{-0.2cm}
	\begin{center}
		\large
		\begin{tabular}{rrl}
		\multirow{4}{*}{\includegraphics[width=1.7cm]{UCM.eps} \hspace{1cm}}&Grado: & Ingeniería Electrónica de Comunicaciones\\
		&Asignatura: & Procesamiento de Señales \\
		&Curso:      & 2025/2026 \\
		\end{tabular}
	\end{center}
	\rule{\textwidth}{0.4pt}
	%
	\vspace{-1.5cm}
	\begin{center}
		\color{blue} \Huge
		\begin{tabular}{c}
	        \textbf{Práctica 8}\\
		\textbf{Señales aleatorias}\\
		\end{tabular}
	\end{center}
	\vspace{-0.5cm}
	\rule{\textwidth}{0.4pt}
}
\maketitle
 \thispagestyle{empty}
\pagestyle{fancy}

\lfoot{\includegraphics[width=0.8cm]{UCM.eps}}
\cfoot{\vspace{-0.6cm} Ingeniería Electrónica de Comunicaciones }
\rfoot{\vspace{-0.6cm} \thepage}
\lhead{Práctica}
\rhead{Práctica 8: Señales aleatorias}
\renewcommand{\headrule}{ \vspace{-0.2cm} \rule{\textwidth}{0.4pt}}
\renewcommand{\footrule}{\vspace{0.5cm} \rule{\textwidth}{0.4pt}\\ \vspace{-0.3cm} \\}

%**************************************************************<- para indicar pregunta de la pr·ctica
\newcounter{ctPreg}
\setcounter{ctPreg}{1}
\newcommand{\itemP}{ \item[{\color{blue}\fbox{\textbf{Tarea \arabic{ctPreg}:}}}] \addtocounter{ctPreg}{1}}  %
\newcommand{\itemB}{ \item[{\color{blue}$\bullet$} ]    }
\newcommand{\itemR}{ \item[{\color{blue}$\triangleright$} ]    }       %
%**************************************************************

\vspace{-8em}

\section*{Objetivos}

\begin{itemize}
	\item Aprender a caracterizar el comportamiento de diferentes señales aleatorias.
	\item Utilizar la correlación cruzada para encontrar un patrón en una señal.
	\item Estimar la respuesta en frecuencia de un sistema desconocido.
\end{itemize}

\section*{Generación y análisis de señales aleatorias}

	\begin{itemize}
		\itemP \textbf{Generación y clasificación de señales aleatorias}\\
		Genera muestras de las siguientes señales aleatorias y analiza su comportamiento:
		\begin{itemize}
			\item Determina, a partir del análisis de las expresiones o de la observación de diferentes muestras de la señal, cuales son markovianas, estacionarias y/o ergódicas. 
			\item En aquellos casos en los que sea posible, representa gráficamente la autocorrelación calculada a partir de una secuencia de la señal generada, y la densidad espectral.
		\end{itemize}

		\begin{itemize}
			\itemR $p(x[n])=\mathcal{N}_{(\mu=5,\sigma=4)}(x[n])$	
			\itemR $p(x[n]|x[n-1])=\mathcal{N}_{(\mu=x[n-1],\sigma=0.5)}(x[n]),~p(x[0])=\delta(x)$
			\itemR $p(x[n]|x[n-1])=\delta(x[n]-x[n-1]),~p(x[0])=\mathcal{N}_{(\mu=0,\sigma=0.25)}(x)$
		\end{itemize}

		\itemP  \textbf{Determinación de patrones en una señal con ruido}\\
		Un sensor de distancia por desplazamiento de fase lanza una señal laser modulada a $f=100~\mathrm{MHz}$ y recibe la
		señal reflejada en el objeto a medir. Para saber la distancia al objeto, se utiliza la expresión $d=\frac{c}{4\pi f}\theta$,
		siendo $c$ es la velocidad de la luz, $f$ la frecuencia de modulación y $\theta$ el desfase. 
		
		Obtén las señales con \color{blue}\texttt{[recibida, enviada] = DatosTarea2()} \color{black}
		determina el desfase entre ambas utilizando la correlación cruzada.

  \newpage
  \itemP  \textbf{Determinación del espectro en frecuencia de un sistema}

  Identifica un sistema desconocido a partir de las relaciones que se establecen entre los espectros en potencia de una señal aleatoria y la misma señal filtrada con un filtro LTI.

  \begin{enumerate}
	  \item \textbf{Genera una señal de identificación}, aleatoria ergódica con espectro plano, que nos permite excitar las frecuencias
	  de interés del sistema a identificar.
  		\begin{lstlisting}[style=Matlab-editor,basicstyle=\scriptsize]
N = 200;
u = randn(N, 1);
			\end{lstlisting}

 	  \item \textbf{Excita el sistema con la señal de identificación.}
 	 Para demostrar el método, utilizamos un sistema cuya función de transferencia es $H(s)=\frac{1}{10 s + 1}$.
  		\begin{lstlisting}[style=Matlab-editor,basicstyle=\scriptsize]
Ts = 2;
fs = 1/Ts;
t = [0:N-1]*Ts;

H = tf(1, [10 1]);
y = lsim(H, u, t);
figure(1); hold on;
	plot(t, u, 'b');
	plot(t, y, 'r');
	xlabel('tiempo (s)')
	title('Excitación - Respuesta')
		  \end{lstlisting}

		\item \textbf{Calcula la densidad espectral cruzada} de ambas señales. En este ejemplo lo hacemos a través de la correlación cruzada con \texttt{xcorr}.
			\begin{lstlisting}[style=Matlab-editor,basicstyle=\scriptsize]
M = floor(N/4);
phi_uu = xcorr(u, M);
phi_uy = xcorr(y, u, M);
Phi_uu = fft(phi_uu);
Phi_uy = fft(phi_uy);

figure(2)
subplot(221);
	plot(abs(Phi_uu));
	title('|\Phi_{uu}(w)|')
subplot(222);
	plot(abs(Phi_uy));
	title('|\Phi_{uy}(w)|')
subplot(223);
	plot(angle(Phi_uu));
	title('arg \Phi_{uu}(w)')
subplot(224);
	plot(angle(Phi_uy));
	title('arg \Phi_{uy}(w)')
		\end{lstlisting}

\newpage
	\item \textbf{Estimar la respuesta en frecuencia} del filtro desconocido, dividiendo
la densidad espectral cruzada entrada/salida entre la densidad espectral de la entrada.
En este ejemplo, representamos el resultado junto con la respuesta en frecuencia real, para
comparar.
		\begin{lstlisting}[style=Matlab-editor,basicstyle=\scriptsize]
K = 2*M + 1;
f = [0:K-1] * fs/K;
w = 2*pi*f;
Hw = Phi_uy ./ Phi_uu;
Hw_ = 1 ./ (10j*w + 1);
figure(3);
subplot(211); hold on;
	plot(f, abs(Hw), 'o');
	plot(f, abs(Hw_), 'r:');
	title('|H(j\omega)|')
	xlim([0, fs/2]);
	ylim([0, 1])
subplot(212); hold on;
	plot(f, angle(Hw));
	plot(f, angle(Hw_));
	xlim([0, fs/2]);
	ylim([-pi, pi])
	title('arg H(j\omega)')
		\end{lstlisting}
  \end{enumerate}

\end{itemize}



\end{document}
